% 基础页面设置
\usepackage{geometry}
\geometry{left=2.8cm,right=2.5cm,top=2.6cm,bottom=2.5cm}

\usepackage{fontspec}
\usepackage{xeCJK}
\usepackage{iftex}
\usepackage{graphicx}
\usepackage{booktabs}
\usepackage{array}
\usepackage{longtable}
\usepackage{multirow}
\usepackage{amsmath,amssymb}
\usepackage{caption}
\usepackage{subcaption}
\usepackage{float}
\usepackage{xcolor}
\usepackage{hyperref}
\usepackage{enumitem}
\usepackage{listings}
\usepackage{fancyhdr}
\usepackage{titlesec}
\usepackage{tocloft}
\usepackage{setspace}
\usepackage{makecell}
\usepackage{tabularx}
\usepackage{chngcntr}
\usepackage{etoolbox}

% 字体：不随包分发字体文件，只使用系统/TeX 发行版中已有字体。
% 推荐 XeLaTeX 编译。若 Windows 有 SimSun/SimHei，会优先使用；否则回退到 Noto/Fandol。
\IfFontExistsTF{Times New Roman}{\setmainfont{Times New Roman}}{\setmainfont{TeX Gyre Termes}}

\IfFontExistsTF{SimSun}{
  \setCJKmainfont[AutoFakeBold=2.5]{SimSun}
  \setCJKsansfont[AutoFakeBold=2.5]{SimHei}
  \setCJKmonofont{FangSong}
  \newCJKfontfamily\songtifont[AutoFakeBold=2.5]{SimSun}
  \newCJKfontfamily\heitifont[AutoFakeBold=2.5]{SimHei}
  \newCJKfontfamily\fangsongfont{FangSong}
}{
  \IfFontExistsTF{Noto Serif CJK SC}{
    \setCJKmainfont[AutoFakeBold=2.5]{Noto Serif CJK SC}
    \setCJKsansfont[AutoFakeBold=2.5]{Noto Sans CJK SC}
    \setCJKmonofont{Noto Sans Mono CJK SC}
    \newCJKfontfamily\songtifont[AutoFakeBold=2.5]{Noto Serif CJK SC}
    \newCJKfontfamily\heitifont[AutoFakeBold=2.5]{Noto Sans CJK SC}
    \newCJKfontfamily\fangsongfont{Noto Serif CJK SC}
  }{
    \setCJKmainfont[AutoFakeBold=2.5]{FandolSong-Regular}
    \setCJKsansfont[AutoFakeBold=2.5]{FandolHei-Regular}
    \setCJKmonofont{FandolFang-Regular}
    \newCJKfontfamily\songtifont[AutoFakeBold=2.5]{FandolSong-Regular}
    \newCJKfontfamily\heitifont[AutoFakeBold=2.5]{FandolHei-Regular}
    \newCJKfontfamily\fangsongfont{FandolFang-Regular}
  }
}

% 正文：宋体，小四；行距建议 1.5 倍。
\AtBeginDocument{\songtifont\zihao{-4}\selectfont\onehalfspacing}

% 标题格式：一级标题黑体二号，二级标题黑体小三，三级标题黑体四号。
\ctexset{
  section = {
    format = \centering\heitifont\zihao{2}\bfseries,
    name = {},
    number = \arabic{section},
    aftername = \quad,
    beforeskip = 1.0\baselineskip,
    afterskip = 0.8\baselineskip
  },
  subsection = {
    format = \heitifont\zihao{-3}\bfseries,
    name = {},
    number = \thesection.\arabic{subsection},
    aftername = \quad,
    beforeskip = 0.8\baselineskip,
    afterskip = 0.5\baselineskip
  },
  subsubsection = {
    format = \heitifont\zihao{4}\bfseries,
    name = {},
    number = \thesubsection.\arabic{subsubsection},
    aftername = \quad
  }
}

% 图表编号按节编号：图 4-1，表 8-2 等。
\counterwithin{figure}{section}
\counterwithin{table}{section}
\renewcommand{\thefigure}{\thesection-\arabic{figure}}
\renewcommand{\thetable}{\thesection-\arabic{table}}
\captionsetup{font={small},labelfont={bf},labelsep=quad}
\renewcommand{\figurename}{图}
\renewcommand{\tablename}{表}

% 页眉页脚
\pagestyle{fancy}
\fancyhf{}
\fancyhead[C]{\zihao{5} 基于 EEG 增强视觉语言模型的退化图像语义识别与生成式描述研究}
\fancyfoot[C]{\zihao{5} \thepage}
\renewcommand{\headrulewidth}{0.4pt}
\renewcommand{\footrulewidth}{0pt}

% 目录格式
\renewcommand{\contentsname}{目\quad 录}
\setcounter{tocdepth}{2}

% 超链接
\hypersetup{
  colorlinks=true,
  linkcolor=black,
  citecolor=black,
  urlcolor=blue,
  pdfauthor={作者姓名},
  pdftitle={基于 EEG 增强视觉语言模型的退化图像语义识别与生成式描述研究}
}

% 列表间距
\setlist{nosep,leftmargin=2em}

% 代码格式
\lstset{
  basicstyle=\ttfamily\small,
  breaklines=true,
  frame=single,
  numbers=left,
  numberstyle=\tiny,
  tabsize=2,
  columns=fullflexible,
  keywordstyle=\bfseries,
  commentstyle=\itshape,
  showstringspaces=false
}

% 常用占位图命令：图片未完成时用占位框，后期替换为自己绘制的图。
\newcommand{\placeholderfigure}[3]{%
\begin{figure}[H]
  \centering
  \fbox{\parbox[c][#3][c]{0.86\linewidth}{\centering\zihao{-4} #2}}
  \caption{#1}
  \label{fig:#2}
\end{figure}
}

% 项目信息命令，可在 cover 中统一修改。
\newcommand{\reporttitle}{基于 EEG 增强视觉语言模型的退化图像语义识别与生成式描述研究}
\newcommand{\college}{计算机与通信工程学院}
\newcommand{\major}{电子信息工程}
\newcommand{\classno}{1901-XX}
\newcommand{\studentid}{2019XXXX}
\newcommand{\studentname}{姓名}
\newcommand{\advisor}{指导教师}
\newcommand{\checkdate}{2026 年 06 月 26 日}
